\chapter{Security Considerations}\label{ch:security}
% Owner: M3 (Shri Krishna P). Draft 1, 12 Nov 2026; sources: blueprint §8, app/core/security.py,
% app/api/deps.py, app/domain/access.py, app/domain/credentials.py, docs/DECISIONS.md.

Every control is enforced on the server; the frontend hides buttons a role cannot use, but
never relies on that. The requirements' six minimum controls map as follows.

\section{Controls}
\begin{longtable}{@{}>{\raggedright\arraybackslash}p{3.2cm}>{\raggedright\arraybackslash}p{12.2cm}@{}}
\toprule
\textbf{Requirement} & \textbf{How I-ITSM meets it} \\
\midrule
\endhead
Authentication & Email and password, or Google Sign-In. A successful sign-in returns a
  signed JWT (HS256) that expires after 30 minutes of clock time; every request after that
  carries it as a bearer token. Expiry is judged by the injected clock, so it is tested
  exactly. The WebSocket authenticates with the same token in its first message and is
  closed (4401) without one or when it expires (D-033). \\
Google Sign-In & The browser obtains an ID token from Google Identity Services; the server
  verifies its RS256 signature against Google's published keys, the audience, issuer and
  expiry, and the \code{hd} claim against the campus domain (D-035). A first sign-in creates an
  Employee; only an administrator grants any other role. \\
Password protection & Argon2 hashes (never stored or logged in clear), 8--128 characters.
  Five failed sign-ins lock the account for 15 minutes. Every failure --- unknown email,
  wrong password, locked or deactivated account --- returns the same answer, so the endpoint
  reveals nothing about which accounts exist; the real reason goes to the audit trail
  (D-019). \\
Role-based authorisation & Five roles map to sets of permissions (\code{Role} $\rightarrow$
  \code{frozenset[Permission]}); code asks \code{user.can(permission)}, never ``is this a
  Technician?''. A FastAPI dependency checks the permission on every route. Nobody can create,
  change or deactivate an account above their own level (D-029). \\
Protection against unauthorised operations & Ownership is checked per incident, in the
  domain (\code{app/domain/access.py}), for every read and action: an employee sees and acts
  only on their own incidents (BR-01); only the assignee, the escalation owner, the team's
  lead or an administrator works one (BR-02); only the reporter verifies or rejects a fix
  (BR-03). A refusal is a 403 and is audited. An integration test calls every endpoint as
  every role and checks the answer against the authorisation matrix. \\
Input validation & Pydantic validates every request body, query and path parameter at the
  boundary (types, lengths, ranges, enums); the domain validates again what only it can know
  (a blank resolution summary, a rating outside 1--5). Errors are 422 with field-level
  messages. Every query goes through the ORM, which binds parameters, so no SQL is built
  from strings. CSV exports neutralise cells that a spreadsheet would run as formulas. \\
Secure API access & In production the API refuses to start with a weak signing secret, a fake
  clock or a CORS policy other than the web app's own \code{https://} origin (D-037); the
  hosts terminate HTTPS. The developer-only clock endpoints do not exist outside the test and
  demo environments (D-009). \\
\bottomrule
\end{longtable}

\section{Integrity and accountability}
\begin{itemize}
  \item \textbf{Hash-chained audit trail.} Each record stores the hash of the one before it:
    \[ h_i = \mathrm{SHA256}\left(h_{i-1} \,\|\, \mathrm{record}_i\right) \]
    Changing, deleting or reordering any record breaks the chain from that point. Appends
    lock the head of the chain so concurrent writers cannot fork it, and the administrator's
    ``verify integrity'' action recomputes it and names the first broken record (BR-10;
    tested by tampering with a row directly in the database).
  \item \textbf{What is audited:} every state change, assignment, escalation, priority
    override (with its reason), AI decision (with its source and confidence), sign-in and
    failed sign-in, and configuration change. Audit records are written in the same
    transaction as the change, so a change cannot exist without its record.
  \item \textbf{Nothing is hard-deleted:} incidents end CLOSED or CANCELLED; users are
    deactivated.
\end{itemize}

\section{Errors that reveal nothing}
All failures are subclasses of one hierarchy, mapped once to HTTP codes; a global handler turns
anything unexpected into a plain 500 and logs it on the server.
\begin{center}
\begin{tabularx}{\textwidth}{@{}lXl@{}}
\toprule
\textbf{Group} & \textbf{Specific errors} & \textbf{HTTP} \\
\midrule
AuthenticationError & InvalidLoginError & 401 \\
AuthorizationError & UnauthorizedOperationError & 403 \\
ValidationError & InvalidUserInputError, InvalidIncidentError & 422 \\
NotFoundError & ServiceNotFoundError, IncidentNotFoundError & 404 \\
ConflictError & InvalidStateTransitionError, DuplicateIncidentError,
  ConcurrentModificationError & 409 \\
AssignmentError & TechnicianUnavailableError, WorkloadExceededError & 409 \\
ConfigurationError & MissingSLAConfigurationError (falls back to the strictest default and
  alerts the admins when a policy is missing) & 503 \\
InfrastructureError & DatabaseFailureError, NetworkFailureError (rollback; the frontend shows
  an offline banner and retries) & 503 \\
\bottomrule
\end{tabularx}
\end{center}
A test sends failing requests of every kind and checks that no response contains a stack trace,
a password hash or a driver's error text.

\section{Secrets and the supply chain}
\begin{itemize}
  \item Settings come from environment variables; \code{.env} is ignored by git and
    \code{.env.example} lists every variable without values. Keys for Gemini, Telegram and
    Google live only in the host's secret store.
  \item gitleaks scans every commit (pre-commit hook) and every pull request (CI).
  \item Adapters never put a key into an error message or a log line (a test checks that a
    Telegram failure does not carry the bot token).
  \item Dependencies are pinned in lock files, and the access token lives in the browser
    tab's \code{sessionStorage} rather than \code{localStorage}, so it ends with the tab (D-022).
\end{itemize}

\section{Known limitations}
\begin{itemize}
  \item There is no rate limiting beyond the per-account lockout; a public deployment should
    add it at the proxy.
  \item Tokens are not stored, so a stolen one stays valid until it expires (30 minutes).
    Every request re-reads the user, though, so a deactivated account, or one whose role has
    changed since the token was issued, is refused at once.
  \item The seeded demo accounts share a published password; they are for evaluation only and
    are to be deactivated afterwards.
\end{itemize}
